%% ============================================================
%%  SECCIÓN 10 — Gráficas
%%  Responsable: Vera Vivanco, Jesús Rodolfo
%% ============================================================

\section{Gráficas}
\label{sec:graficas}

En esta sección se presenta la representación gráfica de los resultados obtenidos para facilitar la comparación visual entre los coeficientes de dilatación lineal medidos y sus respectivos valores de referencia teóricos.

\begin{figure}[H]
    \centering
    \includegraphics[width=0.85\linewidth]{grafica_alfa}
    \caption{Coeficiente de dilatación lineal $\alpha$ determinado experimentalmente para el cobre, aluminio y vidrio borosilicato. Los puntos azules con barras verticales indican el valor central obtenido y su respectiva incertidumbre absoluta propagada $\delta\alpha$. Las marcas rojas en forma de aspa (X) señalan los valores teóricos de referencia bibliográficos aceptados para cada material. Figura generada mediante un script de Python \cite{script_graficas} que utiliza las bibliotecas NumPy \cite{numpy} y Matplotlib \cite{matplotlib} para el procesamiento y graficado de los datos.}
    \label{fig:grafica_alfa}
\end{figure}

La Fig.~\ref{fig:grafica_alfa} evidencia gráficamente la magnitud relativa de los coeficientes calculados y su cercanía respecto a los valores de referencia para cada muestra. El análisis detallado de la dispersión de los puntos, la magnitud de las barras de error y la correspondencia teórica se desarrolla en la Sección~\ref{sec:analisis_graficas}.